\documentclass[11pt]{amsart}
\usepackage{amssymb}
\usepackage{amsmath}
\usepackage{amsthm}
\usepackage{mathrsfs} % \mathscr{}
\usepackage[all, poly]{xy}
\usepackage[x11names,svgnames]{xcolor}
\usepackage{mathtools}
\usepackage[top = 1in, bottom = 1in, left = 1.2in, right=1.2in]{geometry}


\begin{document}

\parskip=0.2in \parindent=0in

\newcommand{\F}{\mathbb{F}}
\newcommand{\C}{\mathbb{C}}
\renewcommand{\Im}{\operatorname{Im}}
\renewcommand{\l}{\overset}
\newcommand{\into}{\hookrightarrow}
\newcommand{\onto}{\twoheadrightarrow}
\newcommand{\tto}{\longrightarrow}
\newcommand{\too}[1]{\overset{#1}\to}
\newcommand{\intoo}[1]{\overset{#1}\into}
\newcommand{\ontoo}[1]{\overset{#1}\onto}
\newcommand{\mapstoo}[1]{\overset{#1}\mapsto}
\newcommand{\bto}{\leftarrow}
\newcommand{\isom}{\cong}
\newcommand{\dsum}{\oplus}
\newcommand{\dsums}{\bigoplus}
\newcommand{\tensor}{\otimes}
\newcommand{\Hom}{\operatorname{Hom}}
\renewcommand{\u}{\underbracket[0.7pt]} % requires mathtools

\newcommand{\inj}{\mathtt{inj}}
\newcommand{\injindex}{\mathtt{inj}\underline{\ }\mathtt{index}}
\newcommand{\inverseind}{\mathtt{inverse}\underline{\ }\mathtt{ind}}
\newcommand{\qut}{\mathtt{qut}}
\newcommand{\phibar}{\bar{\phi}}

\newtheoremstyle{gloss}{\topsep}{\topsep}{}{0pt}{\bfseries}{}{\newline}{\newline *{\bf #3} }
\theoremstyle{gloss}
\newtheorem*{defstar}{Definition}
\newtheorem{theorem}{Theorem}

\newtheoremstyle{newplain}{20pt}{0pt}{\it}{0pt}{\bfseries}{.}{1ex}{}
\theoremstyle{newplain}

\newtheorem{corollary}[theorem]{Corollary}
\newtheorem{claim}[theorem]{Claim}
\newtheorem{lemma}[theorem]{Lemma}
\newtheorem{proposition}[theorem]{Proposition}
\newtheorem{fact}[theorem]{Fact}

\newtheoremstyle{newtextthm}{20pt}{0pt}{}{0pt}{\bfseries}{.}{1ex}{}
\theoremstyle{newtextthm}
\newtheorem{definition}[theorem]{Definition}
\newtheorem{example}[theorem]{Example}
\newtheorem{problem}[theorem]{Problem}
\newtheorem{remark}[theorem]{Remark}
\newtheorem{notation}[theorem]{Notation}


\title{Yoneda products in Guozhen's code}
\date{August 22, 2026}
\maketitle

\textit{I believe everything here goes through if you replace $(\F_2[\tau],
A_\C)$ with $(BP_*, BP_*BP)$.}


Start with a cofree resolution
$$ \F_2[\tau] \to G_0 \to G_1\to G_2 \to \dots $$
of $\F_2[\tau]$ over the dual $\C$-motivic Steenrod algebra $A_\C$.
Suppose we have a cohomology class $\F_2[\tau]\to G_s$, and we want to construct a lift
$$ \xymatrix{
\F_2[\tau]\ar[r]\ar[rd] & G_0\ar[r]\ar[d] & G_1\ar[r]\ar[d] & \dots\ar[r] & G_k\ar[d]^-{\phi_k}\ar[r] &
G_{k+1}\ar[d]^-{\phi_{k+1}}\ar[r] & \dots
\\ & G_s\ar[r] & G_{s+1}\ar[r] & \dots\ar[r] & G_{s+k}\ar[r] & G_{s+k+1}\ar[r] & \dots
}$$
Suppose we have $\phi_k$ and we want to construct $\phi_{k+1}$. We will also
use the fact that the structure maps $G_k\to G_{k+1}$ are constructed by a
composition $G_k\ontoo{\qut} X_{k+1}\intoo{\inj} G_{k+1}$ where $\qut$ is the
quotient map to the cokernel of the previous structure map, and $\inj$ is the
inclusion into a new cofree module.

We want to find an $A_\C$-comodule map $\phi_{k+1}$ that makes the diagram
commute.
\begin{equation}\label{eq:orig-square} \xymatrix{
G_k\ar@{->>}[r]^\qut\ar[d]_-{\phi_k} & X_{k+1}\ar@{^(->}[r]^-{\inj} & G_{k+1}\ar@{.>}[d]^-{\phi_{k+1}}
\\G_{s+k}\ar@{.>>}[r]_\qut & \ar@{^(->}[r]_-{\inj} & G_{s+k+1}
}\end{equation}
The idea is to \textbf{use the universal property of cofree modules}, which
involves mapping \emph{into} a cofree module. Write $G_{s+k+1}= \dsums_{I_{s+k+1}} A_\C$
where $I_{s+k+1}$ is a set of cogenerators.
So we will first construct a
diagram
$$ \xymatrix{
G_k\ar@{->>}[r]^-\qut\ar[d]_-{\phi_k} & X_{k+1}\ar@{^(->}[r]^-{\inj} &
G_{k+1}\ar@{.>}[rd]^-{\phibar_{k+1}}
\\G_{s+k}\ar@{.>>}[r]_-\qut & X_{s+k+1}\ar@{^(->}[r]_-{\inj} &
\u{G_{s+k+1}}_{\dsums_{I_{s+k+1}} A_\C}
\ar[r]_-\epsilon &  \dsums_{I_{s+k+1}} \F_2[\tau]
}$$
I claim it suffices to find $\phibar_{k+1}$.

\begin{lemma}
If $\phibar_{k+1}$ makes the diagram above commute, then \eqref{eq:orig-square}
will commute when $\phi_k$ is defined as the composite
$$ G_k \too{\psi} A_\C \tensor G_k \too{1\tensor \phibar_k} A_\C\tensor
\dsums_{I_{s+k+1}} \F_2[\tau] \too{\text{scalar mult}} \dsums_{I_{s+k+1}}A_\C =
G_{s+k+1}$$
where $\psi$ is the coaction on the cofree comodule
(coaction on each free summand).
\end{lemma}
\begin{proof}
There is a cofree-forgetful adjunction
$$ \Hom_{A_\C\text{-comod}}(M, A_\C\tensor_{\F_2[\tau]} N) \too{\isom}\Hom_{\F_2[\tau]\text{-mod}}(M,N) $$
where the forwards map is post-composition with the counit $\epsilon:A_\C\to
\F_2[\tau]$, and the backwards map is $\phibar \mapsto (1\tensor \phibar)\circ
\psi$. In particular, two $A_\C$-comodule maps $G_k\to G_{s+k+1}$ are the same if and only if
they are the same after post-composition with $\epsilon$.
\end{proof}

In order to construct $\phibar_{k+1}$, we claim that there is a decomposition of
$G_{k+1}$ into free $\F_2[\tau]$-modules,
$$ G_{k+1} \isom \Im(\inj) \dsum R $$
for some $R$. One might worry that there is some element $x\in
G_{k+1}$ that is not in the image, but $\tau x$ is. However, this is OK because
of the construction of $\inj$ and $\qut$.

The important point is that construction of the
resolution comes with a function $\injindex$ that takes a cogenerator
$x$ of $G_{k+1}$ and outputs an element $y\in X_{k+1}$ such that $\inj(y)=u
\cdot x + (\text{higher Curtis terms})$ where $u$ is a unit. {\color{Gray}Isn't 1 the
only unit in $\F_2[\tau]$? Regardless, this is relevant for the odd-primary
generalization.}
In particular, these start with a cogenerator and so are not $\tau$-divisible.
(These terms are the pivot rows in the row-reduced
matrix for $\inj$, and $x$ (or its position) is the pivot index.) Thus we can take:
\begin{itemize} 
\item Basis for $\Im(\inj) = $ these terms $u\cdot x+(\text{higher Curtis terms})$ 
\item Basis for $R = $ a list of non-pivot rows (cogenerators that do not appear
in the image up to higher terms) collected in $\inverseind$.
\end{itemize}

Unfortunately the fact that we don't get a clean $X_{k+1}$-preimage for cogenerators $x\in
G_{k+1}$, only up to higher terms, means that we have to do another cancelation
step. We want to determine $\phibar_{k+1}$ on cogenerators $x$ of $G_{k+1}$.
\begin{itemize} 
\item If $x\in \inverseind$, then define $\phibar_{k+1}(x)=0$. (We just need to
define a $\F_2[\tau]$-linear map, and these are $\F_2[\tau]$-linear generators
of $G_{k+1}$.)
\item Otherwise, we can use $\injindex$ to write $u\cdot x + h = \inj(y)$ for some
$y\in X_{k+1}$ where $u$ is a unit and $h$ represents ``higher terms'' in the Curtis ordered
basis. Then 
$$\phibar_{k+1}(\inj(y)) = u\cdot \phibar_{k+1}(x) + \phibar_{k+1}(h)$$
is determined for us by chasing the diagram: take any preimage along $\qut$ and
then apply $\inj\circ \qut\circ \phi_k$. To solve for
$\phibar_{k+1}(x)$, we use relations like the one above to build a system of
equations over $\F_2[\tau]$ where the pivot positions represent
$\phibar_{k+1}(x)$.

Now row-reduce the system of equations and solve for $\phibar_{k+1}(x)$
using pre-existing row-reduction code (\verb=gaussian=, \verb=symplify_to_led= in
\verb=matrices_mem/4.h=, \verb=matrices/9.h=).
This works because the entries in pivot positions are all units (this is $u$
above).
\end{itemize}

Some potential confusions:
\begin{itemize} 
\item $\mathtt{tauPoly}$ actually consists of $\F_2[\tau]$-monomials only, and
adding $\tau^a + \tau^b$ for $a\neq b$ will silently result in zero.
\item $\inj$ and $\qut$ as matrices only contain $\mathtt{tauPoly}$ entries, but
this is actually correct (there are no hidden higher polynomial terms).
\end{itemize}



%\printbibliography

\end{document}
